% ============================================================
%  EXAMEN TEMA 1: EMPEZANDO A TRABAJAR CON ÁNGULOS
%  Curso 4to Año — Matemática
%  Duración: 90 minutos
%  Temas: Concepto y clasificación de ángulos, sistemas de
%         medición angular y ángulos en el plano cartesiano
%  Autor: Ing. Gabriel Astudillo
%  Nota: enunciado limpio (sin pistas ni desarrollo). Las
%  soluciones están en la "Clave de Respuestas" al final.
% ============================================================

\documentclass[12pt, letterpaper]{article}

% ── Paquetes ─────────────────────────────────────────────────

\usepackage{mathwizards-guia}
\mwguia{Examen — Tema 1: Ángulos}{Ing. Gabriel Astudillo $\cdot$ 4to Año}

\begin{document}
% ============================================================

% ── Portada / Título ─────────────────────────────────────────
\mwportada{Examen del Tema 1}{Empezando a trabajar con ángulos}{Clasificación, Sistemas de Medición y Ángulos en el Plano}

\vspace{4pt}
\begin{cuadroinfo}
  \small
  \textbf{Nombre:} \rule{0.55\linewidth}{0.4pt} \quad \textbf{Fecha:} \rule{0.15\linewidth}{0.4pt} \\[6pt]
  \textbf{Tiempo disponible:} 90 minutos \quad$\bullet$\quad \textbf{Puntaje total:} 100 puntos \\[4pt]
  \textbf{Instrucciones:}
  \begin{enumerate}[leftmargin=*]
    \item Lee cada ejercicio con calma antes de resolverlo.
    \item Muestra \emph{todos} tus pasos; en los ángulos del plano, dibuja si te ayuda.
    \item Simplifica siempre las fracciones al convertir a radianes (ej.: $\frac{120\pi}{180} = \frac{2\pi}{3}$).
    \item No olvides agrupar correctamente grados, minutos y segundos.
    \item No se permite el uso de calculadora.
    \item Escribe con letra clara y revisa tu examen antes de entregarlo.
  \end{enumerate}
\end{cuadroinfo}

\vspace{8pt}

% ============================================================
\section{Concepto, Clasificación y Relaciones (34 puntos)}
% ============================================================

\begin{enumerate}[leftmargin=*]
  \item \textbf{(8 pts)} Clasifica cada ángulo según su medida:
        \[
          a)\ 25^{\circ} \qquad b)\ 90^{\circ} \qquad c)\ 135^{\circ} \qquad d)\ 180^{\circ} \qquad e)\ 360^{\circ}
        \]

  \item \textbf{(8 pts)} Calcula el complemento y el suplemento de:
        \[
          a)\ 38^{\circ} \qquad b)\ 72^{\circ}
        \]

  \item \textbf{(9 pts)} Dos ángulos opuestos por el vértice miden $(4x - 15)^{\circ}$ y $(2x + 35)^{\circ}$. Halla el valor de $x$ y la medida de cada ángulo.

  \item \textbf{(9 pts)} Dos ángulos adyacentes sobre una misma recta miden $(3x + 10)^{\circ}$ y $(x + 30)^{\circ}$. Halla el valor de $x$ y la medida de cada ángulo.
\end{enumerate}

\newpage

% ============================================================
\section{Sistemas de Medición Angular (32 puntos)}
% ============================================================

\begin{enumerate}[leftmargin=*]
  \item \textbf{(8 pts)} Convierte a radianes (simplifica la fracción):
        \[
          a)\ 120^{\circ} \qquad b)\ 315^{\circ} \qquad c)\ 270^{\circ}
        \]

  \item \textbf{(8 pts)} Convierte a grados:
        \[
          a)\ \frac{5\pi}{6}\ \text{rad} \qquad b)\ \frac{7\pi}{4}\ \text{rad} \qquad c)\ \frac{2\pi}{3}\ \text{rad}
        \]

  \item \textbf{(8 pts)} En el sistema sexagesimal:
        \begin{enumerate}[label=\alph*)]
          \item Suma: $42^{\circ} 35' + 27^{\circ} 48'$.
          \item Resta: $90^{\circ} - 37^{\circ} 25'$.
          \item Expresa $48^{\circ} 30'$ en grados con decimales.
        \end{enumerate}

  \item \textbf{(8 pts)} Un ángulo mide $\dfrac{3\pi}{5}\ \text{rad}$. \textquestiondown Cuántos grados mide? Además, expresa $40^{\circ}$ en radianes.
\end{enumerate}

\newpage

% ============================================================
\section{Ángulos en el Plano Cartesiano (34 puntos)}
% ============================================================

\begin{enumerate}[leftmargin=*]
  \item \textbf{(9 pts)} Determina el cuadrante (o el eje) donde se ubica cada ángulo:
        \[
          a)\ 100^{\circ} \qquad b)\ 250^{\circ} \qquad c)\ 310^{\circ} \qquad d)\ 700^{\circ} \qquad e)\ -150^{\circ}
        \]

  \item \textbf{(8 pts)} Escribe dos ángulos coterminales de $210^{\circ}$: uno positivo y uno negativo.

  \item \textbf{(8 pts)} El lado terminal de un ángulo en posición normal pasa por el punto $(-1, \sqrt{3})$. \textquestiondown Cuánto mide el ángulo?

  \item \textbf{(9 pts)} \textquestiondown Qué ángulo forman las manecillas de un reloj a las 4:00? Justifica tu respuesta indicando la posición de cada manecilla.
\end{enumerate}

\newpage

% ============================================================
\ifmwclaves
  \section{Clave de Respuestas}
  % ============================================================

  \subsection*{Sección 1: Concepto, Clasificación y Relaciones}

  \begin{enumerate}[leftmargin=*, label=\textbf{\arabic*.}]
    \item a) agudo ($<90^{\circ}$) \quad b) recto \quad c) obtuso (entre $90^{\circ}$ y $180^{\circ}$) \quad d) llano \quad e) completo.

    \item a) Complemento: $90^{\circ} - 38^{\circ} = 52^{\circ}$; \quad suplemento: $180^{\circ} - 38^{\circ} = 142^{\circ}$.
          \\[2pt]
          b) Complemento: $90^{\circ} - 72^{\circ} = 18^{\circ}$; \quad suplemento: $180^{\circ} - 72^{\circ} = 108^{\circ}$.

    \item Como son opuestos por el vértice, son iguales:
          \[
            4x - 15 = 2x + 35 \Rightarrow 2x = 50 \Rightarrow x = 25.
          \]
          Cada ángulo mide $4(25) - 15 = 85^{\circ}$ (y $2(25) + 35 = 85^{\circ}$).

    \item Por estar sobre una misma recta, son suplementarios:
          \[
            (3x + 10) + (x + 30) = 180 \Rightarrow 4x + 40 = 180 \Rightarrow 4x = 140 \Rightarrow x = 35.
          \]
          Los ángulos miden $3(35) + 10 = 115^{\circ}$ y $35 + 30 = 65^{\circ}$ (suman $180^{\circ}$).
  \end{enumerate}

  \newpage

  \subsection*{Sección 2: Sistemas de Medición Angular}

  \begin{enumerate}[leftmargin=*, label=\textbf{\arabic*.}]
    \item a) $120^{\circ} \cdot \dfrac{\pi}{180^{\circ}} = \dfrac{120\pi}{180} = \dfrac{2\pi}{3}$ rad.
          \quad b) $315^{\circ} \cdot \dfrac{\pi}{180^{\circ}} = \dfrac{315\pi}{180} = \dfrac{7\pi}{4}$ rad.
          \quad c) $270^{\circ} \cdot \dfrac{\pi}{180^{\circ}} = \dfrac{270\pi}{180} = \dfrac{3\pi}{2}$ rad.

    \item a) $\dfrac{5\pi}{6} \cdot \dfrac{180^{\circ}}{\pi} = \dfrac{5(180^{\circ})}{6} = 150^{\circ}$.
          \quad b) $\dfrac{7\pi}{4} \cdot \dfrac{180^{\circ}}{\pi} = 7(45^{\circ}) = 315^{\circ}$.
          \quad c) $\dfrac{2\pi}{3} \cdot \dfrac{180^{\circ}}{\pi} = 2(60^{\circ}) = 120^{\circ}$.

    \item a) $42^{\circ} 35' + 27^{\circ} 48' = 69^{\circ} 83'$. Como $83' = 1^{\circ} 23'$, el resultado es $70^{\circ} 23'$.
          \quad b) $90^{\circ} = 89^{\circ} 60'$; entonces $89^{\circ} 60' - 37^{\circ} 25' = 52^{\circ} 35'$.
          \quad c) $48^{\circ} 30' = 48^{\circ} + \dfrac{30}{60}^{\circ} = 48{,}5^{\circ}$.

    \item $\dfrac{3\pi}{5} \cdot \dfrac{180^{\circ}}{\pi} = \dfrac{3(180^{\circ})}{5} = 108^{\circ}$. \quad
          Y $40^{\circ} \cdot \dfrac{\pi}{180^{\circ}} = \dfrac{40\pi}{180} = \dfrac{2\pi}{9}$ rad.
  \end{enumerate}

  \newpage

  \subsection*{Sección 3: Ángulos en el Plano Cartesiano}

  \begin{enumerate}[leftmargin=*, label=\textbf{\arabic*.}]
    \item a) $100^{\circ}$: cuadrante II. \quad b) $250^{\circ}$: cuadrante III. \quad c) $310^{\circ}$: cuadrante IV.
          \quad d) $700^{\circ} - 360^{\circ} = 340^{\circ}$: cuadrante IV.
          \quad e) $-150^{\circ} + 360^{\circ} = 210^{\circ}$: cuadrante III.

    \item Positivo: $210^{\circ} + 360^{\circ} = 570^{\circ}$. \quad Negativo: $210^{\circ} - 360^{\circ} = -150^{\circ}$.

    \item El punto $(-1, \sqrt{3})$ está en el cuadrante II. El ángulo de referencia con el eje $X$ negativo cumple $\tan \theta = \dfrac{\sqrt{3}}{1} = \sqrt{3}$, es decir, $60^{\circ}$. Por tanto, el ángulo es $180^{\circ} - 60^{\circ} = 120^{\circ}$.

    \item A las 4:00, la manecilla de las horas apunta al 4, que corresponde a $\dfrac{4}{12}(360^{\circ}) = 120^{\circ}$; la manecilla de los minutos apunta al 12 ($0^{\circ}$). El ángulo entre ellas es $|120^{\circ} - 0^{\circ}| = 120^{\circ}$.
  \end{enumerate}

\fi
\end{document}
